\documentclass[../DoAn.tex]{subfiles}
\begin{document}

Phụ lục này em đặc tả thêm hai use case không được trình bày chi tiết trong Chương \ref{chapter:Related_works} do giới hạn dung lượng chương, nhưng vẫn là các chức năng quan trọng của hệ thống NihongoFlow.

\section{Đặc tả use case ``Xem tổng quan tiến độ và streak học tập''}

\begin{table}[H]
\centering
\small
\renewcommand{\arraystretch}{1.3}
\begin{tabular}{|p{2.2cm}|p{1.4cm}|p{2.4cm}|p{7.0cm}|}
\hline
\textbf{Mã usecase} & \multicolumn{1}{l|}{UC06} & \textbf{Tên usecase} & Xem tổng quan tiến độ và streak học tập \\ \hline
\textbf{Tác nhân} & \multicolumn{3}{l|}{Học viên} \\ \hline
\textbf{Tiền điều kiện} & \multicolumn{3}{l|}{Học viên đã đăng nhập.} \\ \hline
\multirow{6}{2.2cm}{\textbf{Luồng sự kiện chính}} & \textbf{STT} & \textbf{Thực hiện bởi} & \textbf{Hành động} \\ \cline{2-4}
 & 1 & Học viên & Truy cập trang \texttt{/dashboard}. \\ \cline{2-4}
 & 2 & Hệ thống & Gọi đồng thời các API: danh sách khóa học đã đăng ký kèm \% hoàn thành (\texttt{completedLessons / totalLessons}), streak hiện tại và dài nhất, danh sách lượt làm quiz gần đây. \\ \cline{2-4}
 & 3 & Hệ thống & Hiển thị thẻ thống kê streak, danh sách khóa học đang học kèm thanh tiến độ, và bảng lịch sử làm quiz (tên bài học, điểm số, ngày làm bài, trạng thái đạt/không đạt). \\ \cline{2-4}
 & 4 & Học viên & Nhấp vào một lượt làm quiz trong lịch sử. \\ \cline{2-4}
 & 5 & Hệ thống & Gọi \texttt{\seqsplit{GET /api/v1/users/me/quiz-attempts/:id}} và mở \texttt{AttemptReviewModal}, hiển thị 10 câu hỏi nhóm theo loại, tô màu xanh cho đáp án đúng và đỏ cho đáp án đã chọn sai. \\ \hline
\multirow{3}{2.2cm}{\textbf{Luồng sự kiện thay thế}} & \textbf{STT} & \textbf{Thực hiện bởi} & \textbf{Hành động} \\ \cline{2-4}
 & A1 & Hệ thống & Học viên chưa đăng ký khóa học nào: danh sách khóa học hiển thị trạng thái rỗng kèm thông báo gợi ý khám phá khóa học. \\ \cline{2-4}
 & A2 & Hệ thống & Chuỗi ngày học bị gián đoạn quá 24 giờ kể từ \texttt{lastActivityDate}: \texttt{currentStreak} được tính lại bằng 0 ở phía server trước khi trả về. \\ \hline
\textbf{Hậu điều kiện} & \multicolumn{3}{l|}{Không thay đổi dữ liệu hệ thống (use case chỉ đọc).} \\ \hline
\end{tabular}
\renewcommand{\arraystretch}{1}
\caption{Đặc tả use case Xem tổng quan tiến độ và streak học tập}
\label{table:B.1}
\end{table}

\section{Đặc tả use case ``Ẩn/hiện khóa học'' (Quản trị viên)}

\begin{table}[H]
\centering
\small
\renewcommand{\arraystretch}{1.3}
\begin{tabular}{|p{2.2cm}|p{1.4cm}|p{2.4cm}|p{7.0cm}|}
\hline
\textbf{Mã usecase} & \multicolumn{1}{l|}{UC07} & \textbf{Tên usecase} & Ẩn/hiện khóa học \\ \hline
\textbf{Tác nhân} & \multicolumn{3}{l|}{Quản trị viên} \\ \hline
\textbf{Tiền điều kiện} & \multicolumn{3}{l|}{Người dùng đã đăng nhập với vai trò \texttt{ADMIN}; khóa học cần thao tác đã tồn tại.} \\ \hline
\multirow{6}{2.2cm}{\textbf{Luồng sự kiện chính}} & \textbf{STT} & \textbf{Thực hiện bởi} & \textbf{Hành động} \\ \cline{2-4}
 & 1 & Quản trị viên & Truy cập trang \texttt{/admin/khoa-hoc}, xem bảng danh sách khóa học kèm cột trạng thái hiển thị. \\ \cline{2-4}
 & 2 & Quản trị viên & Nhấp nút ẩn/hiện tương ứng với một khóa học. \\ \cline{2-4}
 & 3 & Hệ thống & Gọi \texttt{\seqsplit{PATCH /api/v1/admin/courses/:id}} với trường \texttt{hidden} đảo ngược giá trị hiện tại. \\ \cline{2-4}
 & 4 & Hệ thống & Cập nhật bản ghi \texttt{Course} và trả về DTO đã cập nhật. \\ \cline{2-4}
 & 5 & Hệ thống & Cập nhật lại danh sách khóa học (invalidate cache TanStack Query) và hiển thị trạng thái mới trong bảng. \\ \hline
\multirow{3}{2.2cm}{\textbf{Luồng sự kiện thay thế}} & \textbf{STT} & \textbf{Thực hiện bởi} & \textbf{Hành động} \\ \cline{2-4}
 & A1 & Hệ thống & Học viên đã đăng ký khóa học trước khi bị ẩn: vẫn tiếp tục truy cập được các bài học đã đăng ký -- cờ \texttt{hidden} chỉ ảnh hưởng khả năng xuất hiện trong danh sách/khám phá khóa học mới. \\ \cline{2-4}
 & A2 & Hệ thống & Quản trị viên (vai trò \texttt{ADMIN}) luôn nhìn thấy và truy cập được khóa học bất kể trạng thái \texttt{hidden}. \\ \hline
\textbf{Hậu điều kiện} & \multicolumn{3}{p{11.0cm}|}{Trường \texttt{Course.hidden} được đảo trạng thái (\texttt{true} \(\leftrightarrow\) \texttt{false}); nếu \texttt{hidden = true}, khóa học không xuất hiện trong danh sách và không thể truy cập trực tiếp bởi học viên.} \\ \hline
\end{tabular}
\renewcommand{\arraystretch}{1}
\caption{Đặc tả use case Ẩn/hiện khóa học}
\label{table:B.2}
\end{table}

\end{document}
